\section{Related Work}


\textbf{MCP Benchmarks.} Lin et al.~\cite{lin2025large} and Wu et al.~\cite{wu2025mcpzoolargescaledatasetrunnable} develop large-scale maintainable and deployable collections of  MCP servers. 
% MCPCorpus, a large-scale, continuously maintainable collection of real-world MCP servers and clients. 
%  present MCPZoo, a large-scale dataset of deployable MCP servers that can be interacted with. 
Song et al.~\cite{song2025help} introduce MCPGAUGE, a benchmark designed to systematically evaluate how LLMs utilize MCP tools. Wu et al.~\cite{wu2025mcpmark} propose MCPMark, which assesses multi-step workflows with CRUD operations using expert-curated tasks. Yang et al.~\cite{yang2025mcpsecbench} present MCPSecBench, integrating a taxonomy of 17 attack types. Additionally, Zhang et al.~\cite{zhang2025mcp} develop MSB, the first end-to-end evaluation suite measuring LLM agents' robustness against MCP-specific attacks. Fan et al.~\cite{fan2025mcptoolbench++} propose MCPToolBench++, covering a wide range of single- and multi-step interaction tasks. % while addressing prior limitations in dataset coverage, response diversity, and context window constraints.

\textbf{MCP Security.} Song et al.~\cite{song2025beyond} identify novel MCP attack vectors, including tool poisoning, puppet attacks, rug pulls, and exploitation via malicious external resources. Expanding on MCP vulnerabilities, Wang et al.~\cite{wang2025mpma} uncover a new Preference Manipulation Attack. Addressing the need for systematic defense mechanisms, Jing et al.~\cite{jing2025mcip} develop a fine-grained taxonomy and benchmark to detect and mitigate unsafe behaviors in LLM-MCP interactions. Bühler et al.~\cite{buhler2025securing} introduce AgentBound, an access control framework combining declarative, Android-inspired policies with a lightweight enforcement engine. In a similar vein, Kumar et al.~\cite{kumar2025mcp} propose MCP Guardian, enhancing MCP interactions with authentication, rate-limiting, logging, tracing, and WAF scanning. Xing et al.~\cite{xing2025mcp} develop MCP-Guard, a layered defense system for LLM-MCP interactions that integrates static analysis, neural detection, and an LLM-based arbitrator.
